\section{Error Analysis and Discussion}
\label{sec:discussion}

\subsection{Overview of Error Sources}

The experimental results show that pipeline errors are not concentrated in
a single component but may propagate across multiple stages. A retrieval
error may remove a gold event or gold rule; consequently, the gold path is
absent from the candidate set, the verifier evaluates a different path, and
the answer generator produces an answer that is internally consistent with
that incorrect path.

Conversely, even when the gold path is present in the candidate set, the
system may still select the wrong primary path. In this case, retrieval
coverage may be high while Top-1 Exact Path, answer, and citation scores
remain low. The principal error categories are summarized in
Table~\ref{tab:error-summary}.

\begin{table}[t]
    \centering
    \caption{Principal error sources observed in the experiments.}
    \label{tab:error-summary}
    \begin{tabular}{p{0.28\linewidth}p{0.22\linewidth}p{0.38\linewidth}}
        \hline
        \textbf{Error source} & \textbf{Quantitative signal} &
        \textbf{Effect} \\
        \hline
        Path retrieval &
        Error rate \(28.80\%\) &
        The gold path is absent from the candidate set. \\

        Path ranking &
        Error rate \(36.00\%\) &
        The gold path has been retrieved but is not selected as the primary
        path. \\

        Reverse reasoning &
        Answer F1 \(32.67\%\) &
        An incorrect seed or predecessor leads to an incorrect path,
        decision, and answer. \\

        Verification abstention &
        \(20\) supported samples predicted as uncertain &
        Reduces accuracy but avoids issuing an incorrect affirmative
        decision. \\

        Citation selection &
        Citation F1 \(62.86\%\) &
        Omits a gold article or adds an article from a competing path. \\

        Metric for nonlinear structures &
        Exact path is \(0\%\) for branch/convergence &
        A linear metric does not fully reflect answer quality. \\
        \hline
    \end{tabular}
\end{table}

\subsection{The Gap Between Path Retrieval and Path Ranking}

One of the most notable findings is the gap between Top-1 Exact Path and
Oracle Exact Path. Top-1 reaches only \(35.20\%\), whereas Oracle reaches
\(71.20\%\). Thus, for a substantial proportion of questions, the retriever
has generated a candidate path compatible with the gold path, but the
scoring function does not rank that path first. Top-1 is the deployable
path-selection result, whereas Oracle is a diagnostic measure of candidate
set potential rather than a deployable system score.

Three principal factors may explain this phenomenon.

\paragraph{Multiple paths share a prefix or mediator.}

An intermediate event may lead to several different outcomes through rules
from multiple legal articles. For example, events such as ``dangerous
recidivism'' or ``organized crime'' may be linked to multiple sentencing
ranges. These paths share a seed or mediator and therefore receive similar
semantic and graph scores. If the query contains information that
distinguishes the intended outcome but this information is not weighted
strongly enough by the scoring function, the system may select a path that
is structurally valid but leads to the wrong specific consequence.

\paragraph{Outcomes have similar semantic content.}

Many effect events describe imprisonment, fines, or criminal liability but
differ in the magnitude of the penalty or the governing article. An
embedding model may place these events close together in vector space. A
small difference in a numerical value or sentencing range can be legally
material without necessarily producing a sufficiently large embedding
distance.

\paragraph{The path score does not fully incorporate answer intent.}

The retriever runs before Step 5 determines the answer intent. Although the
query is used in semantic retrieval, the scoring function does not impose an
explicit constraint that a reverse question should prioritize a path whose
outcome matches the query and return its predecessor, or that a branch
question should retain multiple paths with a shared prefix. Consequently,
path ranking is driven primarily by relevance and graph scores rather than
by the role that each event must play in the answer.

The Oracle result shows that improving candidate generation alone is
insufficient. A higher-priority direction is to construct a query-aware path
reranker in which query intent is determined before or jointly with
retrieval. Such a reranker could use the following features:

\begin{itemize}
    \item the match between the seed event and the causal part of the query;
    \item the match between the outcome event and the consequence part of
    the query;
    \item the expected role of the event being queried;
    \item the match against the complete rule condition and rule effect;
    \item an article number or penalty magnitude mentioned in the question;
    \item the continuity of the complete path;
    \item event specificity rather than only general semantic similarity.
\end{itemize}

\subsection{Path-Retrieval Errors}

The Path Retrieval Error Rate is \(28.80\%\), meaning that nearly three in
ten samples do not contain the complete gold path in the candidate set.
This error cannot be corrected solely by modifying Step 4 or Step 5.

Potential causes include:

\begin{itemize}
    \item the correct seed event is not in the top-\(k\);
    \item the semantic score is low because the wording of the question
    differs from the event name;
    \item the candidate pool is dominated by many semantically similar
    sentencing events;
    \item graph traversal begins from an event that is topically correct but
    has the wrong role;
    \item the per-seed path limit removes the gold path before reranking;
    \item the candidate rule pool does not contain enough gold rules.
\end{itemize}

Rule Recall@5 and Event Recall@5 reach \(69.93\%\) and \(74.66\%\),
respectively, indicating that the retriever often recovers part of a path
but does not always recover every hop. This observation is consistent with
a Top-1 Edge F1 of \(54.13\%\), which is higher than Exact Path but remains
lower than Oracle Event F1.

Potential retrieval improvements include:

\begin{enumerate}
    \item decomposing the query into representations of the cause, mediator,
    and outcome before seed retrieval;

    \item retrieving condition text and effect text separately;

    \item combining dense scores with sparse lexical matching for article
    numbers, monetary ranges, and sentencing ranges;

    \item increasing seed-event diversity instead of retrieving many events
    that are semantically similar;

    \item performing constrained backward search for reverse reasoning;

    \item applying a rule-aware reranker to complete paths rather than
    ranking individual events independently.
\end{enumerate}

\subsection{Reverse Reasoning as the Principal Weakness}

Reverse reasoning is the lowest-performing group on most metrics:

\begin{itemize}
    \item Rule Recall@5: \(46.67\%\);
    \item Event Recall@5: \(42.22\%\);
    \item Top-1 Exact Path: \(8.89\%\);
    \item Oracle Exact Path: \(26.67\%\);
    \item Verification Accuracy: \(60.00\%\);
    \item Answer Token F1: \(32.67\%\);
    \item Citation F1: \(36.32\%\).
\end{itemize}

In forward reasoning, a query generally states the condition event explicitly
and asks for the outcome. The condition event can be used directly as the
seed, after which the graph is expanded along \texttt{CAUSES} edges. In
reverse reasoning, the query generally begins with an outcome and asks for
its cause. An outcome such as ``incurs criminal liability,'' ``receives a
mitigated sentence,'' or a particular penalty range may have many
predecessors under different legal articles. The inverse problem is
therefore more ambiguous.

Although the retriever is configured with \texttt{direction=both}, allowing
reverse edge traversal is insufficient. The system must also recognize that
the outcome mentioned in the query should be fixed as the endpoint and then
rank predecessors against the complete question content.

Among the \(45\) reverse questions, a Verification Accuracy of \(60\%\)
corresponds to \(27\) correct decisions and \(18\) samples assigned to
\texttt{UNCERTAIN}. The other two \texttt{SUPPORTED} samples assigned to
\texttt{UNCERTAIN} belong to the yes/no positive group, yielding the total
of \(20\) verification errors in the confusion matrix.

The answer generator supports the \texttt{INITIAL\_CAUSE} intent, but it can
only return the first cause on the selected primary path. If the upstream
retriever or verifier selects the wrong path, correctly formatting the
answer intent cannot recover the gold cause. This result shows that
reverse-reasoning errors arise primarily before Step 5.

A direct improvement is to perform constrained backward retrieval:

\begin{equation}
\widehat{p}
=
\arg\max_{p}
S(p \mid q)
\quad
\text{subject to}
\quad
\mathrm{outcome}(p)
=
v_q^{\mathrm{outcome}},
\end{equation}

where \(v_q^{\mathrm{outcome}}\) is the outcome event mapped from the query.
After fixing the outcome, the system can rank paths backward toward the seed
event.

\subsection{Verifier Abstention Behavior}

The confusion matrix shows that the verifier does not directly confuse the
two classes represented in the gold data:

\begin{itemize}
    \item no gold \texttt{SUPPORTED} sample is predicted as
    \texttt{REJECT\_DIRECT\_CLAIM};

    \item no gold \texttt{REJECT\_DIRECT\_CLAIM} sample is predicted as
    \texttt{SUPPORTED};

    \item all \(20\) errors are gold \texttt{SUPPORTED} samples assigned to
    \texttt{UNCERTAIN}.
\end{itemize}

This behavior produces \(100\%\) precision for both \texttt{SUPPORTED} and
\texttt{REJECT\_DIRECT\_CLAIM}, but reduces recall for the
\texttt{SUPPORTED} class to \(91.30\%\).

In a legal system, assigning \texttt{UNCERTAIN} when evidence is insufficient
may be safer than issuing an incorrect affirmative conclusion. However, the
current benchmark contains no gold \texttt{UNCERTAIN} samples, so it cannot
determine whether the abstention threshold is appropriately calibrated.

A future benchmark should contain at least three types of samples:

\begin{enumerate}
    \item supported claims;
    \item rejected claims;
    \item claims for which the available grounds are genuinely
    insufficient.
\end{enumerate}

Such a benchmark would support evaluation of the coverage--risk trade-off
rather than only Accuracy and Macro-F1.

\subsection{Effectiveness of Direct-Claim Detection}

All \(20\) yes/no counterexamples are classified correctly. This is a
positive result for query-aware claim detection: the system recognizes that
the question asks whether a direct relation holds and does not treat a
two-step path as evidence for a direct edge.

However, this result primarily follows from the combination of:

\begin{itemize}
    \item a direct indicator in the query;
    \item the edge structure of the graph;
    \item the distinction between a direct edge and an indirect path.
\end{itemize}

It does not necessarily require a hard intervention. Therefore, the
\(100\%\) accuracy on this \(20\)-sample direct-edge-negative subset should
not be used to claim that LegalSCM has an advantage in counterfactual
reasoning.

\subsection{Answer Quality and Output Length}

Answer Token Recall reaches \(75.40\%\), exceeding the Token Precision of
\(65.34\%\). Answers generally cover the gold content but add:

\begin{itemize}
    \item an introductory phrase such as ``The final consequence is'';
    \item the complete causal chain;
    \item an explanation of the decision;
    \item citation markers \([E_i]\).
\end{itemize}

These additions improve explainability but reduce token precision. A
prediction contains an average of \(48.18\) tokens, whereas a gold answer
contains an average of \(45.07\) tokens.

A Normalized Exact Match of \(0\%\) shows that exact string matching is
unsuitable for explanatory answers. Token F1 and ROUGE-L F1 provide more
informative measures, but they still do not fully evaluate legal
correctness.

One improvement would be to separate the output into two fields:

\begin{itemize}
    \item \texttt{short\_answer}: a concise, direct answer;
    \item \texttt{explanation}: the causal chain and legal grounds.
\end{itemize}

The short answer could then be evaluated using exact match or token F1,
while the explanation could be evaluated separately using path grounding
and citations.

\subsection{Effect of Event Canonicalization}

Some event names in the graph contain multiple phrasings joined by the
\texttt{||} symbol. For example, a node may have the following name:

\begin{quote}
``May be exempted from serving the remainder of the sentence
\texttt{||}
Is exempted from serving the remainder of the sentence.''
\end{quote}

Including the entire canonicalized name in an answer increases the number of
tokens absent from the gold answer and reduces naturalness. It may also
cause semantic retrieval to favor events with more linguistic variants.

The following elements should be separated explicitly:

\begin{itemize}
    \item the canonical event ID;
    \item the alias list;
    \item the preferred display name.
\end{itemize}

The retriever could use all aliases, whereas the answer generator would use
only the preferred display name.

\subsection{Citation Analysis}

Citation F1 reaches \(62.86\%\), which is lower than Answer Token F1. The
two principal error sources are citation misses and extra citations.

\paragraph{Citation miss.}

If the retriever selects the correct event but omits a rule on the gold path,
the answer may still contain most of the gold content while omitting the
corresponding article. This causes answer F1 to exceed citation recall.

\paragraph{Extra citation.}

When several rules support the same causal edge or several paths share an
outcome, Step 5 may retain additional verified evidence that is not part of
the gold citation set. If a citation from this additional evidence appears
in the answer, citation precision decreases.

Deduplication by causal edge has limited repeated citations but has not fully
resolved cases in which multiple articles lead to similar events. One
improvement would be to attach citations to individual assertions in the
answer rather than append all citations at the end.

For example, for a chain explanation:

\[
A \xrightarrow{r_1} B \xrightarrow{r_2} C,
\]

the answer should attach citations as follows:

\begin{quote}
Step 1: \(A\) leads to \(B\) [Article \(a_1\)].
Step 2: \(B\) leads to \(C\) [Article \(a_2\)].
\end{quote}

This presentation clarifies which article supports each hop and facilitates
assertion-level citation-grounding evaluation.

\subsection{Branch and Convergence Structures}

Branch reasoning and convergence reasoning pose a problem for the exact-path
metric. A branch may be represented by a set of paths:

\[
\mathcal{P}_{\mathrm{branch}}
=
\{
A \rightarrow B \rightarrow C_1,\,
A \rightarrow B \rightarrow C_2,\,
\ldots
\},
\]

whereas a convergence structure has the form:

\[
\mathcal{P}_{\mathrm{convergence}}
=
\{
A_1 \rightarrow B \rightarrow C,\,
A_2 \rightarrow B \rightarrow C,\,
\ldots
\}.
\]

No single linear sequence fully represents these structures. Consequently,
Top-1 and Oracle Exact Path may both be \(0\%\) even when the answer
generator correctly enumerates multiple branches or convergent conditions.

The experimental results illustrate this issue clearly:

\begin{itemize}
    \item branch reasoning has a Top-1 Exact Path of \(0\%\) but an Answer
    Token F1 of \(69.33\%\);

    \item convergence reasoning has a Top-1 Exact Path of \(0\%\) but an
    Answer Token F1 of \(84.32\%\).
\end{itemize}

A more appropriate evaluator should use:

\begin{itemize}
    \item Path-set Precision, Recall, and F1;
    \item branch outcome coverage;
    \item convergence input coverage;
    \item subgraph edge F1;
    \item graph edit distance;
    \item assertion-level citation accuracy.
\end{itemize}

Because these metrics are not yet available, path results for branch and
convergence questions are treated only as diagnostic.

\subsection{Effect of Data Composition on Difficulty}

The hard group achieves an Answer Token F1 of \(81.25\%\), exceeding the
\(63.90\%\) of the medium group. Verification Accuracy for the hard group is
also \(100\%\), compared with \(88.95\%\) for the medium group.

This result is inconsistent with the conventional interpretation that hard
samples should be more difficult than medium samples. The reason is that the
difficulty label is affected by question type. The hard group contains many
chain, counterexample, and structurally templated samples with clearly
specified gold answers. The medium group contains many reverse questions,
which the current retriever handles poorly.

Consequently, the current difficulty labels do not constitute an
independently calibrated scale. This study does not use these results to
claim that the system performs better on difficult questions. In future
work, difficulty should be defined using measurable factors such as:

\begin{itemize}
    \item the number of competing candidate paths;
    \item the in-degree and out-degree of the mediator;
    \item similarity among outcomes;
    \item the number of supporting rules;
    \item the number of hops;
    \item lexical overlap between the query and the gold event;
    \item expert agreement.
\end{itemize}

\subsection{Explaining the Null Quality Ablation}

The paired results show that the structural SCM and path ablation achieve
identical values on every quality metric. Several factors plausibly explain
this result.

\paragraph{Retrieval is unchanged.}

The two engines are applied after Step 3. Therefore, the initial rules,
events, and candidate paths are identical.

\paragraph{Most factual paths are simple.}

The current corpus consists primarily of rules with one positive condition
and one positive effect. On these simple paths, fixed-point inference and
reachability generally yield the same supported conclusion.

\paragraph{The counterfactual subset is not sensitive to intervention
semantics.}

The twenty counterfactual samples primarily test direct edges. Both the
structural SCM and path ablation can reject the claim based on the
\(A \rightarrow B \rightarrow C\) structure without producing a difference
in counterfactual outcomes.

\paragraph{Complex alternative mechanisms are absent.}

The benchmark does not yet provide sufficient coverage of cases in which
node deletion and hard intervention yield different results, including:

\begin{itemize}
    \item a mediator with multiple input mechanisms;
    \item an intervention that activates an alternative rule;
    \item an outcome with multiple competing structural mechanisms;
    \item a rule with an exception or negation;
    \item simultaneous interventions on multiple mediators;
    \item a conflict that causes an outcome to become
    \(\mathrm{UNKNOWN}\).
\end{itemize}

\paragraph{Step 5 receives the same decision and primary path.}

When the two verifiers return the same final decision and retain the same
primary path and evidence, the answer generator produces the same answer and
citations. The answer metrics therefore cannot distinguish between the two
engines.

This null result does not establish that the two methods are theoretically
equivalent. A hard intervention replaces the mechanism of a variable,
whereas node deletion changes only graph topology and reachability
\cite{pearl2009causality}. The null result shows only that the current
benchmark does not expose this difference.

\subsection{Trade-off Between Semantics and Runtime}

The structural SCM requires an average of \(3.890\) seconds per sample,
whereas path ablation requires \(1.205\) seconds. The structural SCM is
approximately \(3.23\) times slower per sample; the corresponding batch
wall-clock ratio is approximately \(2.85\).

The additional costs arise from:

\begin{itemize}
    \item initializing or accessing LegalSCM;
    \item constructing the factual context;
    \item fixed-point inference;
    \item rule-coverage checks;
    \item performing an intervention on each mediator;
    \item rerunning inference in the counterfactual world;
    \item generating detailed traces;
    \item computing the node-deletion baseline diagnostic.
\end{itemize}

In return, the structural SCM records information that path ablation does not
provide:

\begin{itemize}
    \item the factual state of the mediator and outcome;
    \item the intervention assignment;
    \item the counterfactual outcome;
    \item whether the outcome changes;
    \item disabled rules;
    \item newly activated rules;
    \item conflicting events;
    \item the number of fixed-point iterations.
\end{itemize}

These fields may be useful for auditability and error analysis. However, the
current study does not include an expert evaluation that quantifies the
value of these traces. Their interpretability benefit should therefore not
be recast as an accuracy claim. In particular, every measured quality metric
is identical between the structural SCM and path ablation in the current
paired experiment.

\subsection{Requirements for a Future Counterfactual Benchmark}

To evaluate LegalSCM appropriately, a future benchmark must provide ground
truth not only for the final decision but also for world states and
intervention traces. A counterfactual sample should include:

\begin{itemize}
    \item the factual context;
    \item the intervened mediator;
    \item the factual mediator state;
    \item the factual outcome;
    \item the intervention assignment;
    \item the counterfactual outcome;
    \item an outcome-changed flag;
    \item disabled mechanisms;
    \item alternative mechanisms;
    \item an expert-validated explanation.
\end{itemize}

The benchmark should also contain paired cases in which node deletion and
hard intervention produce different results. Only then can a paired
ablation test the added value of structural semantics.

\subsection{Implications for Practical Deployment}

The current results show that the pipeline can produce grounded answers for
the entire benchmark and correctly classify the direct-edge negative claims.
However, a Top-1 Exact Path of \(35.20\%\), a Citation F1 of \(62.86\%\),
and weak reverse-reasoning performance are not sufficient to support
deployment as an automated legal advisory system. The \(35.20\%\) Top-1
score is the deployable path-selection result; the \(71.20\%\) Oracle score
is diagnostic and must not be presented as deployed performance.

In a practical application, the output should be presented as an
evidence-grounded suggestion, and users should be able to:

\begin{itemize}
    \item inspect the source rules and articles;
    \item verify each hop in the causal path;
    \item inspect alternative supporting paths;
    \item recognize when the system returns \texttt{UNCERTAIN};
    \item distinguish factual inference from counterfactual intervention;
    \item identify the corpus version and update date.
\end{itemize}

The system should neither conceal uncertainty nor display only the final
answer without enabling evidence inspection.

\subsection{Improvement Priorities}

Based on the experimental results, improvements are prioritized as follows:

\begin{enumerate}
    \item \textbf{Query-aware path reranking.}
    Incorporate answer intent and event role into scoring before selecting
    the primary path.

    \item \textbf{Backward retrieval for reverse reasoning.}
    Fix the outcome mentioned in the query and retrieve conditional
    predecessors.

    \item \textbf{Hybrid lexical--dense retrieval.}
    Add sparse matching for article numbers, monetary amounts, and
    sentencing ranges.

    \item \textbf{Preferred event display name.}
    Separate the canonical ID, aliases, and display name to remove
    \texttt{||} strings from answers.

    \item \textbf{Assertion-level citation.}
    Attach each article to a specific hop or assertion.

    \item \textbf{Subgraph evaluation.}
    Add path-set metrics for branch and convergence structures.

    \item \textbf{Intervention-sensitive benchmark.}
    Construct expert-validated gold factual and counterfactual worlds.

    \item \textbf{Extend LegalSCM.}
    Support multiple conditions, negation, exceptions, conflicts, and rule
    priorities.
\end{enumerate}

These improvements should be implemented in a new iteration and evaluated
through a new paired experiment. They must not be mixed with the frozen
v4.1/v5.1 results in the current study.
